% Prose rendering of PrimeTensor/Tensor/Calculus.lean.
% Fragment: no \documentclass, no preamble, no document environment.

\section{Abstract multiplicative differential calculus}
\label{Tensor:Calculus:sec}

\leanfile{PrimeTensor/Tensor/Calculus.lean}

\subsection*{What this file establishes}

Nothing, in the sense of proving: the file contains no theorem.  It fixes
vocabulary --- points, scalar fields, vector fields --- and defines the three
vector-calculus operators the fluid layer will need, gradient, divergence and
Laplacian, on top of an abstract carrier of directional derivatives.

The module header states the reason for the abstraction: it deliberately
separates the differential \emph{identities} from the eventual analytic
construction of a derivative on completed prime barcodes.  A
$\name{Differential}$ here is a bare family of operators with no axioms, so
every identity proved about these three operators downstream will be an
identity about folds and components, provable before any analysis exists and
true of whatever concrete derivative is eventually supplied.

What is worth reading closely is that the operators are \emph{products}, not
sums --- Design note~\ref{Tensor:Calculus:note:product} --- and that the folds
require no unit, which is a consequence of the decision in
\texttt{PrimeTensor/Depth.lean} that nothing starts at zero ---
Remark~\ref{Tensor:Calculus:rem:no-unit}.

\subsection{Fields}

\begin{definition}[\lean{Point}, \lean{ScalarField} and \lean{VectorField}]
\label{Tensor:Calculus:def:fields}
For a type $X$ of coordinate values, a type $A$ of field values and a dimension
$d$,
\begin{align*}
  \name{Point}(X,d) &\defeq \name{Axis}(d) \to X, \\
  \name{ScalarField}(X,A,d) &\defeq \name{Point}(X,d) \to A, \\
  \name{VectorField}(X,A,d) &\defeq
    \name{Point}(X,d) \to \name{Tensor.Vector}(A,d).
\end{align*}
\end{definition}

\begin{leancode}
/-- A point in a positive-dimensional coordinate space. -/
abbrev Point (X : Type) (dim : Depth) := Axis dim → X

/-- Scalar field. -/
abbrev ScalarField (X A : Type) (dim : Depth) := Point X dim → A

/-- Vector field. -/
abbrev VectorField (X A : Type) (dim : Depth) := Point X dim → Tensor.Vector A dim
\end{leancode}

All three are $\ctor{abbrev}$s, so they unfold on sight and carry no wrapping.
The coordinate type $X$ receives no structure anywhere in the file: a point is
a tuple of values indexed by an axis, and nothing may be done with it but feed
it to a field.  By \texttt{PrimeTensor/Depth.lean} the axis is nonempty, so a
point always has at least one coordinate.

\subsection{The carrier of derivatives}

\begin{definition}[\lean{Differential}]\label{Tensor:Calculus:def:differential}
$\name{Differential}(X,A,d)$ is the structure with the single field
\[
  \name{d} : \name{Axis}(d) \to
    \name{ScalarField}(X,A,d) \to \name{ScalarField}(X,A,d),
\]
assigning to each coordinate direction an operator on scalar fields.
\end{definition}

\begin{leancode}
structure Differential (X A : Type) (dim : Depth) where
  d : Axis dim → ScalarField X A dim → ScalarField X A dim
\end{leancode}

\begin{remark}[No axioms whatever]\label{Tensor:Calculus:rem:no-axioms}
The structure has one field and no laws.  A $\name{Differential}$ is
\emph{any} assignment of an endomorphism of the scalar-field space to each
axis: nothing requires it to be a morphism of the value carrier, to satisfy a
product rule, to commute with itself across directions, or to bear any relation
to $\name{HasMulDerivativeAt}$ of
\texttt{PrimeTensor/Analysis/Derivative.lean}.  The two notions of derivative
in the development are, at this point, entirely unconnected, and the header
says that is deliberate.

The consequence is a division of labour.  Everything provable about
Definition~\ref{Tensor:Calculus:def:gradient} and the two below is provable by
unfolding folds and components, with no analysis and no hypotheses --- and so
holds for whichever concrete $\name{Differential}$ is eventually constructed.
The price is that nothing provable here has analytic content.
\end{remark}

\subsection{The three operators}

\begin{definition}[\lean{Differential.gradient}]\label{Tensor:Calculus:def:gradient}
$\name{gradient}(D,f)$ is the vector field whose $i$-th component at $x$ is
$D.\name{d}\,i\,f\,x$.
\end{definition}

The gradient imposes no structure on $A$ at all --- it only repackages the
family of directional derivatives as a rank-one tensor in the sense of
\texttt{PrimeTensor/Tensor/Basic.lean}.

\begin{definition}[\lean{Differential.divergence} and \lean{laplacian}]
\label{Tensor:Calculus:def:div-lap}
For $A$ carrying a multiplication,
\begin{align*}
  \name{divergence}(D,u)(x)
    &\defeq \name{fold}\bigl((\cdot\,\cdot),\, d,\;
       i \mapsto D.\name{d}\,i\,(y \mapsto (u\,y).\name{component}\,i)\,x\bigr), \\
  \name{laplacian}(D,f)(x)
    &\defeq \name{fold}\bigl((\cdot\,\cdot),\, d,\;
       i \mapsto D.\name{d}\,i\,(D.\name{d}\,i\,f)\,x\bigr).
\end{align*}
\end{definition}

\begin{leancode}
/-- Multiplicative divergence: product of diagonal directional derivatives. -/
def divergence {X A : Type} [Mul A] {dim : Depth}
    (D : Differential X A dim) (u : VectorField X A dim) : ScalarField X A dim :=
  fun x =>
    Axis.fold (· * ·) dim (fun i =>
      D.d i (fun y => (u y).component i) x)
\end{leancode}

\begin{designnote}[Products where the classical operators have sums]
\label{Tensor:Calculus:note:product}
Both operators fold with multiplication.  Classically
\[
  \operatorname{div} u = \sum_i \partial_i u_i,
  \qquad
  \Delta f = \sum_i \partial_i \partial_i f,
\]
and the sums here are products.  This is the same substitution the whole
development makes, and it reads correctly in the logarithmic chart of
\texttt{PrimeTensor/Analysis/Derivative.lean}: since $\log$ turns a product
into a sum, the multiplicative divergence of $u$ is the ordinary divergence of
the log-transformed field, and likewise for the Laplacian.  Nothing is lost and
nothing additive is introduced.

Two features of the classical operators survive unchanged.  The divergence
contracts a vector field to a scalar by pairing the $i$-th direction with the
$i$-th component --- the diagonal, as the docstring says.  And the Laplacian
takes only \emph{same-axis} second derivatives, $D_i D_i f$: no mixed partial
$D_i D_j f$ with $i \neq j$ appears in this file, and no commutation of mixed
partials is available to be assumed, since
Remark~\ref{Tensor:Calculus:rem:no-axioms} leaves the operators unrelated.
\end{designnote}

\begin{remark}[The fold needs no unit, and that is the point of positive depth]
\label{Tensor:Calculus:rem:no-unit}
The two definitions require $[\ctor{Mul}\,A]$ and nothing else.  In particular
they do not require $[\ctor{One}\,A]$, even though a product over an index set
normally needs a unit for the empty case.

There is no empty case.  $\name{fold}$ of \texttt{PrimeTensor/Depth.lean} was
defined with no identity argument precisely because $\name{Axis}(d)$ always has
at least one coordinate, $\name{Depth}$ having no zeroth constructor.  So the
divergence of a vector field over a one-dimensional axis is its single
directional derivative, not that derivative multiplied by an empty product.

This is the concrete payoff of the founding decision, and it matters more here
than a convenience: the eventual value carrier is $\name{MulReal}$, which does
have a unit, but the definition is stated for an arbitrary type with a
multiplication, and that generality is available only because no unit is
needed.
\end{remark}

\begin{remark}[Order and bracketing are not shown to be irrelevant]
\label{Tensor:Calculus:rem:order}
$\name{fold}$ associates to the right and visits the axis in the order
$\ctor{first}$, then the successors, and the definitions inherit that.  Since
$A$ carries only $[\ctor{Mul}\,A]$, with neither associativity nor
commutativity assumed, the value of a divergence genuinely depends on the
bracketing and on the enumeration of the axis.  For a carrier satisfying
$\name{IsMulCarrier}$ of \texttt{PrimeTensor/Structure.lean} both dependencies
vanish, but this file neither assumes that nor proves the independence.  Any
later identity that reorders a fold will have to establish it.
\end{remark}

\begin{remark}[What is absent]\label{Tensor:Calculus:rem:absent}
There is no curl, no second-order mixed operator, and no theorem of any kind
--- no linearity, no product rule, no identity relating divergence to gradient,
and no analogue of $\operatorname{div}\operatorname{grad} = \Delta$, which in
this notation would say that the divergence of the gradient of $f$ is the
Laplacian of $f$ and is in fact true by unfolding, the two folds having the
same summand.  It is not stated.  Nor is any $\name{Differential}$ constructed:
the structure has no instance anywhere in this file, so the three operators are
so far defined and never applied.
\end{remark}

\subsection*{Summary of the correspondence}

\begin{center}
\small
\begin{tabular}{@{}lll@{}}
\hline
Lean declaration & Kind & Here \\
\hline
\lean{Point}                      & abbrev     & Def.~\ref{Tensor:Calculus:def:fields} \\
\lean{ScalarField}                & abbrev     & Def.~\ref{Tensor:Calculus:def:fields} \\
\lean{VectorField}                & abbrev     & Def.~\ref{Tensor:Calculus:def:fields} \\
\lean{Differential}               & structure  & Def.~\ref{Tensor:Calculus:def:differential} \\
\lean{Differential.gradient}      & definition & Def.~\ref{Tensor:Calculus:def:gradient} \\
\lean{Differential.divergence}    & definition & Def.~\ref{Tensor:Calculus:def:div-lap} \\
\lean{Differential.laplacian}     & definition & Def.~\ref{Tensor:Calculus:def:div-lap} \\
\hline
\end{tabular}
\end{center}

\noindent
Every declaration of \texttt{PrimeTensor/Tensor/Calculus.lean} appears above,
together with the single field $\name{d}$ of $\name{Differential}$.  The file
contains no theorem, and so no \lean{sorry}, no axiom, and no unproved named
proposition.  Design note~\ref{Tensor:Calculus:note:product} and
Remarks~\ref{Tensor:Calculus:rem:no-axioms}, \ref{Tensor:Calculus:rem:no-unit},
\ref{Tensor:Calculus:rem:order} and~\ref{Tensor:Calculus:rem:absent} are stated
for the reader; the logarithmic reading of the two folds is a gloss and is not
formalised.
